\section{High-jet differences of the actual quantum step}
\label{cdq-section}

The background contraction can be combined with the actual
common-domain fluctuation integral. The result compares two
admissible complete inherited inputs. It does not assert that
the quantum step preserves their common low jets.

Use the graphs, radius $b$, exponent $a_0=1/16$ and parameters
of Theorem~\ref{cdm-theorem}, with the additional choices below.
Pad absent coefficients by zero so that both inputs have one
common connected-carrier index set. Write
\begin{equation}
 E_j=\sum_HE_{j,H},\qquad
 \norm{E_j}_{a_0;b}\le\bar\epsilon,\qquad
 D=E_1-E_2,\qquad D_H=E_{1,H}-E_{2,H}.
 \label{cdq-inputs}
\end{equation}
Both total inputs satisfy the local gauge and logarithm-branch
hypotheses of \eqref{cdm-input}. Assume the full coefficientwise
condition
\begin{equation}
 D^nD_H(0)=0\qquad(H\text{ arbitrary},\ 0\le n\le4),
 \label{cdq-common-jets}
\end{equation}
including every mixed field derivative. Equality of low jets
only after summing carriers is not sufficient.

\subsection{Compatible coefficient families}
\label{cdq-compatibility-section}

Choose one pre-gauge reserve large enough for the stronger
joint exponent $B=53/8$, and use its same supported
$a_p=\sum_Xa_{p,i;X}$ in Sections~\ref{cdj-section} and
\ref{cdh-section}. Fix one sufficiently large $J$, one
$\ell_{\mathrm{tail}}$, the same support truncation
$X\subset B_{\ell_{\mathrm{tail}}}(i)$, the same halo
$K_0(H)=N_{\infty,2}(\pi_RH)$, and the same deterministic
joining trees and carrier grouping. The stronger joint
budgets imply the weaker background budgets.
Also impose the final-radius condition
\begin{equation}
 L_H=\frac{R}{8C_A}\ge2,\qquad bL_H\le r_A.
 \label{cdq-high-disk}
\end{equation}

In the actual fluctuation family each coefficient has the
form $a_{p,i;X}+d_{\lambda;i;X}$ with
$d_{0;i;X}=0$ termwise. Consequently the core/tail expansion
at fluctuation parameter zero is exactly the background
expansion at field-scaling parameter one:
\begin{equation}
 U^{\mathrm{fluct}}_{K,0}[E]
   =U^{\mathrm{field}}_{K,1}[E]
   =(\mathcal T_{R,J}E)_K.
 \label{cdq-background-identity}
\end{equation}
The local coefficient functions, multiindices, support
occurrences and trees coincide before summing $K$.
The two proofs may estimate this same assignment in different
weights. Extra sites of the halo can be retained as declared
support even when the background coefficient does not read them.

Core/tail differentiation of $E_H$ and support grouping are
linear in the input family. Thus $\mathcal T_{R,J}$ and the
inherited joint difference $Q_E$ are linear, and
\begin{equation}
 Q_{E_1}-Q_{E_2}=Q_D,\qquad
 \norm{Q_E}_B\le a_Q\norm{E}_{a_0;b},\qquad
 a_Q=16C_TC_R\frac{gT}{b}.
 \label{cdq-Q-linear}
\end{equation}
This is the actual estimate \eqref{cdm-inherited-bound}.

For the Gaussian logarithm retain all possible carrier
labels, including coefficients that happen to vanish.
Use common forest and hull assignments, precision pieces
and the normally convergent cluster-log branch. Let
$\mathcal L(Q)$ denote this one analytic coefficient-valued
logarithm map with the fixed base exponent and precision.
The compatible actual output family is
\begin{equation}
 \mathscr F(E)=
 \mathcal T_{R,J}E+\Delta+Q_{\det}+\mathcal L(Q_E).
 \label{cdq-output-map}
\end{equation}
Its scalar sum is precisely \eqref{cdm-update}, with the
same germ or gauge-orbit density interpretation. The
classical background, covariance, cutoff, base factors
and external determinant families are identical for
the two inputs. No $E$-dependent quantum saddle is used.

\subsection{A two-input Lipschitz estimate}
\label{cdq-theorem-section}

Set
\begin{equation}
 C_H=C_0(8C_A)^5,\qquad
 C_{\mathrm{quant}}=1024eS_1C_TC_R,
 \label{cdq-constants}
\end{equation}
where $C_0$ is \eqref{cdh-uniform-transport}. It is fixed
before $R$; $C_{\mathrm{quant}}$ is fixed after the finite
geometry is selected. Strengthen the allowed budgets to
\begin{equation}
 v\le\frac{\eta_*}{64e},\qquad
 a_Q\bar\epsilon\le\frac{\eta_*}{128e},\qquad
 e^3\tau(T)\le\eta_*/4,\qquad \eta_*\le1.
 \label{cdq-budgets}
\end{equation}
All precision and domain conditions of
Theorem~\ref{cdm-theorem} remain in force.

\begin{theorem}[Quantum differences with common low jets]
\label{cdq-theorem}
For the compatible families above, condition
\eqref{cdq-common-jets} implies
\begin{equation}
 \norm{\mathscr F(E_1)-\mathscr F(E_2)}_{a_0;b}
 \le
 \left(\frac{C_H}{R}
             +C_{\mathrm{quant}}\frac{gT}{b}\right)
                       \norm{E_1-E_2}_{a_0;b}.
 \label{cdq-full-bound}
\end{equation}
Writing $\mathcal R_{\ge5}$ for the coefficientwise
total-degree remainder in \eqref{cdh-projection-bound},
\begin{equation}
 \norm{\mathcal R_{\ge5}
       [\mathscr F(E_1)-\mathscr F(E_2)]}_{a_0;b}
 \le6\left(\frac{C_H}{R}
             +C_{\mathrm{quant}}\frac{gT}{b}\right)
                       \norm{E_1-E_2}_{a_0;b}.
 \label{cdq-projected-bound}
\end{equation}
These estimates are uniform in old depth $M$ and compatible
volume. They assert a bound for the specified compatible
canonical families, not for arbitrary localizations of
the same summed functional.
\end{theorem}

\begin{proof}
By \eqref{cdq-background-identity}, the actual background
difference is the family to which
Theorem~\ref{cdh-contraction} applies. The coefficientwise
jet condition gives
\begin{equation}
 \norm{\mathcal T_{R,J}D}_{a_0;b}
                  \le\frac{C_H}{R}\norm{D}_{a_0;b}.
 \label{cdq-background-bound}
\end{equation}
The external $\Delta,Q_{\det}$ cancel coefficientwise.
Linearity in \eqref{cdq-Q-linear} leaves
\[
 \mathscr F(E_1)-\mathscr F(E_2)
   =\mathcal T_{R,J}D
       +\mathcal L(Q_{E_2}+Q_D)-\mathcal L(Q_{E_2}).
\]

If $D=0$ the claim is immediate. Otherwise put
\begin{equation}
 q_D=a_Q\norm{D}_{a_0;b},\qquad
 L_D=\frac{\eta_*}{32e q_D}.
 \label{cdq-interpolation-radius}
\end{equation}
Since $\norm{D}\le2\bar\epsilon$,
\eqref{cdq-budgets} gives $q_D\le\eta_*/(64e)$ and
$L_D\ge2$. Interpolate the inherited exponent as
$Q_{E_2}+\zeta Q_D$ on $|\zeta|<L_D$. Its total exponent
with the fixed base satisfies
\begin{equation}
 \norm{V_{\mathrm{base}}+Q_{E_2}+\zeta Q_D}_B
 \le
 \left(\frac1{64}+\frac1{128}+\frac1{32}\right)
              \frac{\eta_*}{e}
 =\frac{7\eta_*}{128e}<\frac{\eta_*}{8e}.
 \label{cdq-exponent-budget}
\end{equation}
Equal-carrier exponentiation costs at most $e^xx$, where
$x$ is this bound. With $\eta_*\le1$ its multiplier norm
is at most $\eta_*/8$. The tail quarter budget then puts
the connected-log norm at most $S_1\eta_*$ throughout the
disk, by \eqref{cdm-log-bound}. Precision and base remain
fixed, and the cluster theorem supplies one common
nonvanishing logarithm branch.

The interpolated full input $E_2+\zeta D$ need not stay
inside the original $\bar\epsilon$ class on this larger
disk. We have instead verified the actual Gaussian
interaction budgets directly. Normal convergence gives
a holomorphic anchored log family. Banach Cauchy, with
the contour integral before the positive anchored sum,
therefore yields
\begin{align}
 \norm{\mathcal L(Q_{E_1})-\mathcal L(Q_{E_2})}_{a_G;b}
 &\le\frac{S_1\eta_*}{L_D-1}
 \le64eS_1q_D\notag\\
 &\le C_{\mathrm{quant}}\frac{gT}{b}\norm{D}_{a_0;b},
 \qquad a_G=\frac12.
 \label{cdq-log-bound}
\end{align}
Both endpoints are evaluations of the same canonical
analytic cluster family, so no independent logarithm
branch or carrier choice is introduced. Since
$a_G\ge a_0$, this proves \eqref{cdq-full-bound}.
Finally \eqref{cdh-projection-bound} has norm at most
six on the unchanged circular product ball and preserves
the declared carriers. Applying it proves
\eqref{cdq-projected-bound}.
\end{proof}

\subsection{A concrete one-half factor and its limitation}
\label{cdq-scope-section}

The parameter order can give a concrete projected factor.
Fix the stronger common reserve and its $C_A,\kappa_0,r_A$,
then choose a finite power $R=\ell^s$ satisfying all
previous compression and amplitude conditions and
\begin{equation}
 R\ge24C_H.
 \label{cdq-R-choice}
\end{equation}
Choose the common $J$ and tail cutoff afterward, then
one final $b$ satisfying \eqref{cdq-high-disk} and all
classical and Gaussian neighborhoods. Finally choose
$g,T(g)$ to meet \eqref{cdq-budgets} and
\begin{equation}
 C_{\mathrm{quant}}gT/b\le1/24.
 \label{cdq-g-choice}
\end{equation}
With a logarithmically growing cutoff the latter holds
for all sufficiently small $g$ at these fixed parameters.
Equations \eqref{cdq-full-bound}--\eqref{cdq-projected-bound}
then give
\begin{equation}
 \norm{\mathscr F(E_1)-\mathscr F(E_2)}_{a_0;b}
       \le\frac1{12}\norm{D}_{a_0;b},\qquad
 \norm{\mathcal R_{\ge5}
       [\mathscr F(E_1)-\mathscr F(E_2)]}_{a_0;b}
       \le\frac12\norm{D}_{a_0;b}.
 \label{cdq-half-contraction}
\end{equation}
The same $b$ is used at input and output. The canonical
map above uses a common field-independent normalization.
Separate output vacuum subtraction may double the
unprojected estimate, but leaves the projected bound
unchanged because $\mathcal R_{\ge5}$ already kills
constants.

The scalar sums compared here are differences of the
actual common-domain quantum map. The proof first
compares two admissible complete inputs; it does not
apply a rooted/pre-gauge identity separately to a
nongauge-invariant extraction. The representation-sensitive
condition \eqref{cdq-common-jets} and the compatible
family construction are essential to the norm statement.

The full output difference can regenerate field jets
of degrees zero through four through the fluctuation
integral. Thus the quantum map does not preserve the
affine input class with the prescribed low jets.
Projecting the output is bounded but need not preserve
gauge invariance, because gauge transformations are
nonlinear in these potential coordinates. No
gauge-invariant low-jet extraction or closed
relevant/marginal flow has been constructed.
This contractive response is therefore not yet a
Banach fixed-point iteration on admissible physical
activities. Large fields, mixed histories, global
measure consistency, extraction and coupling flow,
and continuum observables remain separate.
